\scanpage{28}
\uline{Prop. 7} Soit $A\in\underline b$, alors
\margin{$\Big|$}
\[
\mu_{|A|}=\mu_{|A^*|}qquad\text{i.e.}\qquad
\varphi_{|A|}=\varphi_{|A^*|}.
\]

On a $\varphi_{|A|}^{2}=\varphi_{A^*A}$, $\varphi_{|A^*|}^{2}=\varphi_{A^*A}$, il suffit de prouver $\varphi_{A^*A}=\varphi_{AA^*}$ \inserted{i.e. $\mu_{AA^*}=\mu_{A^*A}$} \deleted{par suite}. Par prop. 2, cor. 4, on peut se borner au cas $A\in\underline a$, mais alors \deleted{il suffit de p\unsure{rouver}} $\mu_{AA^*}$ et $\mu_{A^*A}$ sont deux mesures \deleted{finies} \inserted{à support compact} sur \deleted{$[0,+\infty[$} \inserted{$R^+$}, il suffit de prouver que
\[
\mu_{AA^*}(x^p)=\mu_{A^*A}(x^p)
\]
pour les entiers $p>0$, i.e.
\[
\operatorname{Tr}((AA^*)^p)=\operatorname{Tr}((A^*A)^p),
\]
ce qui résulte immédiatement de $T(RS)=T(SR)$.

$\varphi_{|A|}=\varphi_{|A^*|}$ et les quantités dérivées
$\Phi_{|A|}=\Phi_{|A^*|}$, $\Psi_{|A|}=\Psi_{|A^*|}$ etc \inserted{doivent être utilis\unsure{és|ées}} \deleted{\unsure{mesurent la}} \deleted{une mesure} comme une mesure de l'ordre de grandeur de $A$.

\fbox{N° 6}\quad\uline{Compléments sur les déterminants}

\deleted{\uline{Définition}} On pose \deleted{pour \unsure{tout} $A\in\widetilde{\underline b}$} d'abord, pour $A\in\widetilde{\underline b}$
\[
(1)\qquad\Delta(A)=\deleted{\log}\inserted{\exp}\int_0^{T(1)}\log\varphi_{|A|}.
\]
\deleted{Dans \unsure{ce qui suit}}
\deleted{En d'autres t\unsure{ermes},} \deleted{si $A\in\underline b$} \deleted{$\Delta(A)=\Delta(\unsure{|A|})$}
\deleted{$\Delta(A)=\Delta(|A|)$, et si $t$}

a) Si $T$ est finie. $\Delta(A)$ est définie pour tout $A\in\mathcal A$, et $0\leq\Delta(A)\unsure{<}+\infty$. On a
\[
\Delta(\lambda A)=|\lambda|^{T(1)}\Delta(A),
\qquad\deleted{\Delta(1)=1\quad\unsure{\Delta(A)|\text{et }\Delta(A)}}
\]

b) Si $T$ est infinie, supposons $A\geq0$, $A=\lambda1+B$ ($B\in\underline\Phi$, $B$ hermitien). \deleted{\unsure{Si}}
\mspage{29}
